% -----------------------------------------------------------------------
% §2 Related Work
% -----------------------------------------------------------------------
\section{Related Work}
\label{sec:related}

\subsection{Post-Quantum Key Exchange and Authentication}
\label{sec:related:pqake}

The standardisation of post-quantum cryptographic primitives has spurred a
growing body of work on constructing efficient and provably secure authenticated
key exchange (AKE) protocols from lattice-based building blocks.
CRYSTALS-Kyber~\cite{EuroSP:BDKLLSSSS18}, the KEM underlying NIST FIPS~203
(ML-KEM)~\cite{NIST:FIPS203}, provides IND-CCA2-secure encapsulation whose
security reduces to the hardness of the Module-Learning-With-Errors (MLWE)
problem.
CRYSTALS-Dilithium~\cite{TCHES:DKLLSSS18}, standardised as ML-DSA (FIPS~204)~\cite{NIST:FIPS204},
provides EUF-CMA-secure signatures whose security rests on the Module-Short-Integer-Solution
(MSIS) and MLWE assumptions.
Both schemes offer three security levels ($\{2, 3, 5\}$) aligned with NIST
security categories II, III, and V, giving protocol designers a tunable
security--overhead trade-off.

NIST has also standardised Falcon (FIPS~206, FN-DSA)~\cite{NIST:FIPS206} as a
signature scheme that achieves significantly smaller signatures than ML-DSA: at
the 128-bit security level, a Falcon signature is approximately 4.8 times
smaller than an ML-DSA-65 signature, though its key-generation procedure
requires Gaussian sampling.
For applications that can tolerate larger signatures in exchange for a
different security foundation, SLH-DSA (FIPS~205)~\cite{NIST:FIPS205} provides
a hash-based, stateless signature alternative that does not rely on lattice
assumptions, with signature sizes ranging from 8 to 50~KB.

KEMTLS~\cite{CCS:SchSteWig20} is a landmark engineering contribution that
replaces the signature-based handshake of TLS~1.3 with KEM-based mutual
authentication, eliminating handshake signatures from the critical path entirely.
PQ-Aliro takes a complementary but distinct approach: rather than redesigning
the handshake, it preserves the Aliro APDU command structure and substitutes
only the underlying primitives, ensuring wire-level structural compatibility
with existing access-control infrastructure.
Crucially, KEMTLS targets TLS record-layer semantics over TCP/IP, where
payload fragmentation is handled transparently by the transport layer; PQ-Aliro
must instead work within the hard 255-byte ISO-DEP APDU boundary of the NFC
physical layer, making byte-level communication overhead a primary design
constraint rather than an afterthought.

On the theoretical side, Fujioka et al.~\cite{ASIACCS:FSXY13} construct
one-pass AKE protocols from one-way CCA-secure KEMs, establishing a principled
framework for building AKE security from KEM security in the standard model.
PQ-Aliro's authentication structure follows a similar spirit: the ML-KEM
encapsulation provides session-key establishment, while ML-DSA certificates
bind ephemeral keys to long-term identities.
Related lattice-based AKE constructions with tight security reductions, such as
the work of Pan et al.~\cite{C:PanWagZen23}, confirm that provably secure,
tightly-reduced lattice AKE is achievable, reinforcing the theoretical
foundations on which PQ-Aliro rests.
Earlier work~\cite{PKC:FSXY12} demonstrated strongly secure AKE from lattice
assumptions in the PKC setting, providing further evidence of the maturity of
this research direction.
These theoretical results, however, do not directly address the byte-level
constraints that motivate the design of PQ-Aliro, which we examine in the
context of AKE security models next.

\subsection{AKE Security Models}
\label{sec:related:models}

The security of authenticated key exchange protocols is formalised through a
hierarchy of increasingly powerful adversarial models.
The foundational Bellare--Rogaway (BR93) model~\cite{C:BelRog93} establishes
the core AKE framework: it defines session partnering via matching
conversations, formalises adversarial queries including long-term-key
corruption, and captures session-key indistinguishability and entity
authentication.
Canetti and Krawczyk~\cite{EC:CanKra01} extended this framework with the CK
model, which additionally captures leakage of session-specific internal state
and allows for more refined notions of forward secrecy.
The eCK model of LaMacchia, Lauter, and Mityagin~\cite{PROVSEC:LaMLauMit07}
pushes adversarial power further, simultaneously exposing long-term keys and
ephemeral randomness to the adversary, and is widely regarded as one of the
strongest standard AKE security definitions.

PQ-Aliro's security analysis is cast in a BR-style model augmented with an
identity-privacy experiment.
The privacy dimension is motivated by the work of Lyu et al.~\cite{AC:LLHG22},
who formalise privacy-preserving AKE in the standard model and prove that user
identities can be kept unlinkable across sessions even against an active
adversary.
In the NFC access-control setting, tracking resilience is a practical
requirement: a passive eavesdropper recording NFC transactions must not be able
to link repeated accesses to a single credential holder.
Our identity-privacy experiment, which tests whether an adversary can
distinguish between two target users based on observed APDU transcripts,
follows the definitional approach of~\cite{AC:LLHG22} adapted to the two-party
door-access setting.
The TLS~1.3 cryptographic analysis of Dowling et al.~\cite{JC:DFGS21} provides
a further point of comparison for multi-stage AKE models, although its
record-layer focus is orthogonal to our APDU-layer concerns.
This emphasis on APDU-layer constraints naturally leads to the question of how
PQC primitives perform on the constrained platforms typical of NFC access
control.

\subsection{PQC on Constrained Platforms}
\label{sec:related:constrained}

Deploying post-quantum cryptography in resource- or bandwidth-constrained
environments remains an active area of research.
Lima et al.~\cite{NIST3PQC:Lima21} evaluate the performance of Kyber KEM in a
native Android mobile application, reporting that encapsulation and
decapsulation are fast enough for interactive use on commodity smartphones.
Their findings provide direct empirical support for the feasibility of
mobile-side PQC in PQ-Aliro's Android HCE implementation.
Post-quantum WireGuard~\cite{SP:HNSWZ21} integrates lattice-based KEMs into the
WireGuard VPN protocol, demonstrating that the handshake overhead of
NIST-round-3 KEMs can be absorbed by a constrained network protocol without
breaking functional compatibility.

The above results focus on IP-based transports or general mobile computing
contexts.
Despite these results, a significant gap remains.
Existing work addresses PQC integration in TLS and VPN-class protocols, where
message fragmentation is handled by TCP or UDP at the transport layer and
byte-level overhead does not impose hard per-message limits.
The NFC context is fundamentally different: ISO/IEC~14443 ISO-DEP APDUs carry
at most 255 bytes of application data, and the 424~kbit/s physical-layer
bit rate turns each kilobyte of overhead into measurable latency.
To our knowledge, \emph{no prior work provides a complete, APDU-granular
communication-overhead analysis for a PQC-upgraded NFC access-control
protocol}.

On the standards side, neither Aliro~1.0~\cite{CSA:Aliro10} nor CCC Digital
Key~3.0~\cite{CCC:DigKey3}---the two leading specifications for smartphone-based
digital access control---provides a post-quantum variant, a PQC migration
roadmap, or a formal security proof for their AKE core.
PQ-Aliro addresses all three dimensions---wire compatibility, byte-accurate
overhead characterisation, and formal security proof---simultaneously: it
delivers a wire-compatible PQC upgrade, a byte-accurate overhead
characterisation across all nine ML-KEM/ML-DSA parameter combinations, and a
formal proof of mutual authentication, forward secrecy, and identity privacy.